% =============================================================================
% Lecture 18: Engineering Physical AI
% Physical AI: Machine Learning Systems That Sense and Act (Volume IV)
% =============================================================================
% The course's closing lecture. It spans the whole stack: the four laws join
% into one design argument (segment 1), the argument transfers to a machine the
% course never studied, the packaging cell (segment 2), it runs out at premises
% no sensor reports (segment 3), and the open premises choose the next
% experiment (segment 4). The synthesis returns to Lecture 0's aisle run with
% the finished stopping budget and the full record chain.
%
% Running case: the warehouse mobile manipulator on its aisle run at 1.3 m/s,
% now with a clear hydraulic film across the aisle (Napkin Math 18.2).
%
% Build: make deck V=4 D=18_frontier   (figures: images/make_figs.py)
% Numbers: images/calcs.py prints every derived value on these slides.
% Plan: lecture_map.yml in this folder; Volume IV course plan, Lecture 18.
% =============================================================================

\documentclass[aspectratio=169, 12pt]{beamer}
\usepackage{../../assets/beamerthememlsys}
\mlsyssetup{
  volume={Volume IV},
  chapter={Chapter 18},
  chaptertitle={Engineering Physical AI},
  logo={../../assets/img/logo-mlsysbook.png},
  instlogo={../../assets/img/logo-harvard.png},
}
\usepackage[T1]{fontenc}
\usepackage[scaled=0.9]{helvet}
\usepackage{courier}
\renewcommand{\familydefault}{\sfdefault}
\usepackage{booktabs,amsmath}
\graphicspath{{images/}}
\setcounter{mlsystotalsections}{5}

\title{Engineering Physical AI}
\author{Vijay Janapa Reddi}
\institute{Harvard University}
\date{}

% A large answer choice for predict and concept-question slides
\newcommand{\choice}[2]{\textbf{\textcolor{mlsysaccent}{#1}}\enspace #2}
\newcommand{\rchoice}[2]{\colorbox{mlsysaccent!15}{\choice{#1}{#2}}}
% A small gray tag that marks an illustrative, derived, or measured value
\newcommand{\smalltag}[1]{{\scriptsize\textcolor{midgray}{\textit{#1}}}}
% One line with a text-free thumbnail of the figure that taught it
\newcommand{\thumbline}[2]{%
  \parbox[c]{2.1cm}{\centering\includegraphics[width=2.0cm,height=1.0cm,keepaspectratio]{#1}}%
  \hspace{0.25cm}\parbox[c]{0.76\textwidth}{\raggedright #2}}
% Volume IV roles (volumes/vol4/palette.md)
\definecolor{delaycol}{HTML}{D97706}
\definecolor{plantcol}{HTML}{334155}
\definecolor{nervecol}{HTML}{059669}
\definecolor{conseq}{HTML}{DC2626}
\definecolor{brainnavy}{HTML}{1F407A}

% One cell of the missing-evidence grid: a figure that shows the gap, its class, one line
\newcommand{\gapcell}[3]{%
  \begin{minipage}[c]{0.53\linewidth}\centering\includegraphics[width=\linewidth,height=2.1cm,keepaspectratio]{#1}\end{minipage}%
  \hspace{0.15cm}\begin{minipage}[c]{0.44\linewidth}\raggedright\textbf{#2}\\[0.05cm]#3\end{minipage}}
\begin{document}

% =============================================================================
% TITLE AND HOOK
% =============================================================================
\mlsysnote{[0.5 min] This is the last lecture of the course. Read the thesis aloud once. The lecture adds three ideas of its own: the four-law audit of a new machine, the residual-claims register, and choosing the next experiment by the evidence a premise lacks. Everything else is the course's own material, used at once. The cover shows a layered machine under glass beside a gap labeled unobserved reality; that gap is today's subject. The notes total 71.5 minutes, 58.5 of them [CORE]; cutting every [OPTIONAL] note leaves a coherent lecture of under an hour. FLEX: [CORE].}
\mlsystitle{Engineering Physical AI}{The weakest open premise sets what a machine may be trusted to do, and retiring that premise chooses what to build next.}{cover_frontier_blueprint.png}

\begin{frame}{Predict: does the stop fit on a clear film?}
  \mlsysactivity{Predict}{3 min}
  \begin{columns}[c]
    \begin{column}{0.44\textwidth}
      \centering\safeimg[width=\textwidth,height=4.6cm,keepaspectratio]{pallet-jack.jpg}\\[0.05cm]
      {\scriptsize A hydraulic hand pallet jack on polished concrete}
    \end{column}
    \begin{column}{0.53\textwidth}
      \small A pallet jack's leaking seal leaves a clear film across the aisle. The base runs at 1.3 m/s toward a rack end 1.10 m away. On dry concrete its stopping budget uses 969.0 mm.\\[0.2cm]
      \normalsize How much does it use on the film?\\[0.15cm]
      \choice{A}{969.0 mm}\qquad\choice{B}{1.3 m}\\[0.1cm]
      \alt<2->{\rchoice{C}{2.1 m}}{\choice{C}{2.1 m}}\qquad\choice{D}{no stop at all}\\[0.2cm]
      \uncover<2->{\small Every deadline was met. No sensor on the machine reports the film.}
    \end{column}
  \end{columns}
  \mlsyscredit{Photo: Auledas, CC BY 4.0, Wikimedia Commons. Numbers: Napkin Math 18.2 (floor friction illustrative; 969.0 mm derived, known components only).}
  \note{[3 min] LINK: Lecture 12 chose the 1.3 m/s aisle speed against the 1.10 m clear distance at the rack end. NARRATE: Step 1. The photo is an ordinary hydraulic pallet jack; its lift cylinder runs on oil, and a worn seal leaves a thin, clear film on polished concrete. Read the card, then ask for a vote with fingers, one to four, a minute with a neighbor, and a second vote. Step 2 reveals C: the stopping budget grows from 969.0 mm to about 2.13 m, nearly twice the clear distance. The base hits the rack end. ENGAGE: ask a student who chose A what they assumed. Expected: that the brakes and the timing decide the stop. WARN: students will look for a software fault. There is none. The permission path met every deadline, the chunk lease expired on schedule, and the drives applied braking torque on time. The film changed the friction under the tires, and the cameras and lidar cannot tell clear oil from dry concrete. The rest of the lecture explains why a better model cannot fix this and what can. FLEX: [CORE].}
\end{frame}

\begin{frame}{Which law did each part of the course add?}
  \mlsysactivity{Retrieval}{2.5 min}
  \begin{columns}[c]
    \begin{column}{0.55\textwidth}
      \centering\mlsyskitfig[width=\textwidth,height=5.0cm,keepaspectratio]{stack_roadmap}
    \end{column}
    \begin{column}{0.42\textwidth}
      \only<1|handout:0>{\normalsize Question of the day\\[0.2cm] Write the four laws in order, and one record each part added.\\[0.3cm]{\small\textcolor{midgray}{Two minutes.}}}%
      \only<2-|handout:1>{\small
        \textbf{1 Irreversibility}\\ {\footnotesize limit record (L2)}\\[0.15cm]
        \textbf{2 Endogenous data}\\ {\footnotesize provenance record (L5)}\\[0.15cm]
        \textbf{3 Proposal is not permission}\\ {\footnotesize enforcement record (L12)}\\[0.15cm]
        \textbf{4 Permission must be earned}\\ {\footnotesize release manifest (L16)}}
    \end{column}
  \end{columns}
  \note{[2.5 min] NARRATE: Step 1. The course stack with every lecture placed on its level; Lectures 14 to 17 sit on the governance envelope. Two minutes, from memory: the four laws in order, and one record from each part. Walk the room; a blank line tells you what to review. Step 2: the answer. Lecture 1 named the four laws, and each part's first lecture stated its own: Lectures 2, 5, 8, and 14. Only the first is a physical law; the other three are design principles the book raises to laws by choice. Any record from the right part counts: limit or handoff for Part I; provenance, policy manifest, or evaluation record for Part II; any of the six Part III records; authority, fault, release manifest, or approval record for Part IV. Lecture 17 added the approval record, which names who answers for each decision behind an action. Today adds the sixteenth and last record. FLEX: [CORE].}
\end{frame}

\begin{frame}{Why does a better model not make a better machine?}
  \small
  By the end of today, you will be able to:\\[0.25cm]
  \thumbline{thumb_closed_loop_step2.pdf}{Join the four laws into one design argument for a machine.}\\[0.2cm]
  \thumbline{thumb_audit_step4.pdf}{Run the four-law audit on a machine the course never studied.}\\[0.2cm]
  \thumbline{thumb_register.pdf}{Write a residual-claims entry for a premise no sensor reports.}\\[0.2cm]
  \thumbline{thumb_gaps.pdf}{Choose the next experiment by the kind of evidence a premise lacks.}\\[0.25cm]
  {\footnotesize\textcolor{midgray}{Thread: all four laws. Record: the residual-claims register.}}
  \note{[1.5 min] The driving question is the one the book opened with. NARRATE: read the four goals. Define the key phrase. A premise is a condition a claim depends on, such as the floor's friction under the stopping budget. An open premise is one no evidence yet supports. The record this lecture adds is the residual-claims register, which closes the record chain that Lecture 1 introduced. FLEX: [CORE].}
\end{frame}

% =============================================================================
% SECTION: ONE ARGUMENT
% =============================================================================
\mlsysnote{[0.5 min] Segment 1 assembles the course into one picture and one table. Read the question aloud. FLEX: [CORE].}
\mlsysdivider{How do eighteen lectures join into one argument?}{Engineering Physical AI, \S18.1}
\section{One Argument}

\begin{frame}{One causal loop closes through the world at three rates}
  \begin{columns}[c]
    \begin{column}{0.62\textwidth}
      \centering\mlsyssteps[width=\textwidth,height=5.3cm,keepaspectratio]{closed_loop_step}{2}
    \end{column}
    \begin{column}{0.35\textwidth}
      \centering\safeimg[width=\textwidth,height=3.0cm,keepaspectratio]{erflex-mobile-manipulator.jpg}\\[0.05cm]
      {\scriptsize A mobile manipulator of the running machine's class}\\[0.25cm]
      \uncover<2->{\small A film under the tires reaches no level.}
    \end{column}
  \end{columns}
  \mlsyscredit{Slide variant of \S18.1, figure 1. Photo: ER-FLEX (MiR250 base, UR16e arm), Auledas, CC BY 4.0, Wikimedia Commons.}
  \note{[2.5 min] LINK: Lecture 1's five levels and its proposal boundary. NARRATE: Step 1. The running machine, assembled. The Brain proposes: the intent model at 2 Hz, the chunk policy renewed every 50 ms, the two-speed brain of Lecture 3. Below the proposal boundary, the permission path on the safety microcontroller of Lecture 4 admits, corrects, or refuses on a 1 kHz tick. The drives close current loops at 20 kHz. Four orders of magnitude separate the fastest loop from the slowest. Step 2: the loop closes only through the World, and the World returns only what the sensors capture. ENGAGE: which level would see the oil film? Expected: none; friction under a clear film reaches no sensor, so it is invisible to every level at once. WARN: students treat governance as a sixth level. It is the envelope around all five; it changes on the timescale of releases and sets the budgets the fast loop enforces. FLEX: [CORE].}
\end{frame}

\begin{frame}{Each law leaves records that the next law reads}
  \footnotesize
  \renewcommand{\arraystretch}{1.02}
  \begin{tabular}{@{}P{0.22\textwidth}P{0.38\textwidth}P{0.34\textwidth}@{}}
    \toprule
    \textbf{Law (part)} & \textbf{Records that carry it} & \textbf{Condition to keep true} \\
    \midrule
    \textcolor{plantcol}{\textbf{Irreversibility (I)}} & Limit records; handoff record & Each measured limit, such as floor friction, still holds \\
    \uncover<2->{\textcolor{nervecol}{\textbf{Endogenous data (II)}}} & \uncover<2->{Provenance, policy manifest, evaluation record} & \uncover<2->{The states the machine enters lie inside the evaluated ODD} \\
    \uncover<3->{\textcolor{nervecol}{\textbf{Proposal is not permission (III)}}} & \uncover<3->{Observation contract, belief, intent lease, trajectory, enforcement, placement} & \uncover<3->{Fresh state reaches the permission path, which meets its deadline} \\
    \uncover<4->{\textcolor{brainnavy}{\textbf{Permission must be earned (IV)}}} & \uncover<4->{Authority, fault, release manifest, approval, \textbf{residual-claims register}} & \uncover<4->{The premises of the release claim still hold after release} \\
    \bottomrule
  \end{tabular}
  \note{[2.5 min] NARRATE: one row per step, each the chapter's table 1. Step 1: Part I turns the body's limits into budgets; its records are the limit records of Lecture 2 and the handoff record of Lecture 4. Step 2: Part II confines competence to the states the closed loop visits; the provenance record (Lecture 5), policy manifest (Lecture 6), and evaluation record (Lecture 7). The ODD is the operational design domain of Lecture 5. Step 3: Part III places an independent permission path between proposal and actuator. Step 4: Part IV grants only what the evidence supports, and the residual-claims register, today's record, holds what the evidence does not yet support. ENGAGE: which condition did the oil film break? Expected: row one, a measured limit; the friction behind the credible deceleration of Lecture 2 no longer held. WARN: students read the table as four separate checklists. Each row's condition is an input to the next. A speed limit in row one changes the states row two must evaluate. FLEX: [CORE].}
\end{frame}

\begin{frame}{Guarantees a person once supplied now sit inside the machine}
  \centering
  \begin{minipage}[b]{0.30\textwidth}\centering
    \safeimg[width=\textwidth,height=4.0cm,keepaspectratio]{shakey-1969.jpg}\\[0.05cm]
    {\scriptsize Shakey, 1969: a reflex halt beneath the planner}
  \end{minipage}\hfill
  \begin{minipage}[b]{0.30\textwidth}\centering
    \safeimg[width=\textwidth,height=4.0cm,keepaspectratio]{atlas-drc-door.jpg}\\[0.05cm]
    {\scriptsize DARPA Robotics Challenge, 2015: operators near the loop}
  \end{minipage}\hfill
  \begin{minipage}[b]{0.34\textwidth}\centering
    \safeimg[width=\textwidth,height=4.0cm,keepaspectratio]{erflex-mobile-manipulator.jpg}\\[0.05cm]
    {\scriptsize Today: a permission path and timed handover}
  \end{minipage}
  \mlsyscredit{Photos, Wikimedia Commons: Shakey, The wub, CC BY-SA 4.0; Team ViGIR Atlas at the DRC, DARPA, public domain; ER-FLEX, Auledas, CC BY 4.0. History: \S18.1.}
  \note{[1.5 min] NARRATE: the chapter closes the history Lecture 1 began. Left: SRI's Shakey, whose reflex halt kept working beneath a slow symbolic planner. Middle: an Atlas humanoid at the 2015 DARPA Robotics Challenge door task, where remote operators stayed close to the loop. Right: today's machine. Learning now sits in pooled data such as Lecture 5's Open X-Embodiment, 2.42 million episodes from 27 robot types. The reflex halt became the discrete permission path of Lecture 12; the watchful operator became the intervention authority of Lecture 14, with timed handover budgets; the oversight a person once supplied became the standing conditions of release of Lecture 16. ENGAGE: which guarantee would vanish if a team removed the permission path to save latency? Expected: the reflex halt; the proposer would command the drives directly. FLEX: [OPTIONAL].}
\end{frame}


\begin{frame}{Does an intervention mean the task failed?}
  \mlsysactivity{Vote}{2 min}
  \begin{columns}[c]
    \begin{column}{0.50\textwidth}
      \begin{mlsyscard}{mlsysaccent}
        \small The base enters an inspected route rated for 1 m/s. The governor trims its 1.3 m/s proposal to 1 m/s, and the delivery completes.\\[0.1cm]
        \textbf{True or false: the run is a failure.}
      \end{mlsyscard}
    \end{column}
    \begin{column}{0.46\textwidth}
      \uncover<2->{\normalsize\textbf{False.} The leg took \textbf{30\%} longer.\\[0.25cm]
      \small Three channels, recorded apart:\\[0.05cm]
      \textcolor{brainnavy}{proposer autonomy}\\
      \textcolor{plantcol}{whole-machine task success}\\
      \textcolor{nervecol}{protective intervention}}
    \end{column}
  \end{columns}
  \mlsyscredit{Engineering Physical AI, \S18.5, first fallacy; 30 m leg from \S18.4.2; 30\% derived (images/calcs.py).}
  \note{[2 min] The chapter's first fallacy as a vote. ENGAGE: hands up for true, then for false. Step 2: false. The governor applied a least-change speed correction, the minimal-intervention filter of Lecture 12, and the delivery completed. A 30 m leg took 30 s against 23.1 s, 30 percent longer (derived). WARN: teams that log every intervention as a failure corrupt their own telemetry. Lecture 7 separated the outcome channels: how often the learned model ran without override, whether the machine finished the task in time, and how often protection acted. Mixing them hides whether the bottleneck is distribution shift, perception latency, or a governor tuned too conservatively. Land segment 1. The four laws are one argument, and each part's records are the evidence the next part reads. FLEX: [CORE].}
\end{frame}

% =============================================================================
% SECTION: A NEW MACHINE
% =============================================================================
\mlsysnote{[0.5 min] Segment 2 tests whether the method transfers. The packaging cell is a machine the course has not studied. FLEX: [CORE].}
\mlsysdivider{Does the method carry to a machine the course never studied?}{Engineering Physical AI, \S18.2}
\section{A New Machine}

\begin{frame}{A packaging cell holds all three machine classes}
  \begin{columns}[c]
    \begin{column}{0.40\textwidth}
      \centering\safeimg[width=\textwidth,height=4.2cm,keepaspectratio]{carton-conveyor.jpg}\\[0.05cm]
      {\scriptsize An automated carton conveyor line}
    \end{column}
    \begin{column}{0.58\textwidth}
      \centering\safeimg[width=\textwidth,height=4.2cm,keepaspectratio]{cell_schematic.pdf}
    \end{column}
  \end{columns}
  \mlsyscredit{Photo: TGW Mechanics GmbH, CC BY 3.0, Wikimedia Commons. Schematic: slide variant of \S18.2.2, figure 2a (hypothetical cell).}
  \note{[2 min] LINK: Lecture 1's three machine classes, named for the margin each spends. NARRATE: the photo is a real automated carton line; the chapter's cell is a hypothetical one of this kind, laid out tightly for cycle rate. A carriage carries one carton along a track. An insertion head presses the flaps to seat the retention tabs. A gate closes across the track when the next station is full, and a carriage still moving drives its carton into the gate. ENGAGE: name the class of each part. Expected: the carriage is moving mass (Class 1), the insertion head is contact force (Class 2), and the drive, cycling all shift, is heat (Class 3). Each has an irreversible outcome: a crushed carton and bent frame, a buckled carton wall, and winding insulation that ages and does not recover when it cools. FLEX: [CORE].}
\end{frame}

\begin{frame}{The audit asks one question per law, in order}
  \centering
  \mlsyssteps[width=\textwidth,height=6.1cm,keepaspectratio]{audit_step}{4}
  \mlsyscredit{Engineering Physical AI, \S18.2.1.}
  \note{[2 min] The first concept the lecture owns: the four-law audit of a new machine. NARRATE: one step per law. Step 1: what can the task do that no later command can undo, and which measured bound excludes it? The answer is the limit records, such as a speed ceiling from the stopping budget. Step 2: do the demonstrations and the closed-loop evaluation cover the states inside that envelope? Step 3: which independent path admits, corrects, or refuses each proposal, on which sensors, within which timing budget? Step 4: which claims does the evidence support, under which standing conditions, who answers for each decision, and which premises stay open? WARN: the order matters. Evaluating a policy before the stopping budget fixes the speed ceiling wastes the evaluation on states the permission path will never admit. What transfers is the questions, the budgets, and the test protocols; the warehouse machine's numbers do not. FLEX: [CORE].}
\end{frame}

\begin{frame}{Estimate: can the carriage stop before the gate?}
  \mlsysactivity{Estimate}{3.5 min}
  \begin{columns}[T]
    \begin{column}{0.40\textwidth}
      \begin{mlsyscard}{mlsysaccent}
        \small\textbf{Given} \smalltag{(illustrative)}\\[0.05cm]
        Speed 1.50 m/s\\
        Stop line 60 mm away; keep 15 mm reserve\\
        Credible braking 15 m/s$^2$\\
        Closure to full braking 64 ms\\[0.05cm]
        $d = v\,\tau + v^2/2a$
      \end{mlsyscard}
    \end{column}
    \begin{column}{0.56\textwidth}
      \small
      \uncover<2->{Delay: $1.50 \times 0.064$ = \textbf{96 mm}\\
      Braking: $1.50^2 / 30$ = \textbf{75 mm}\\
      Stop \textbf{171 mm}: it passes the line before braking starts.\\[0.2cm]}
      \uncover<3->{Ceiling on 45 mm:\\
      $v = -a\tau + \sqrt{(a\tau)^2 + 2ad} = 1.51 - 0.96$\\[0.05cm]
      {\Large\textbf{\textcolor{mlsysaccent}{0.55 m/s}}}}
    \end{column}
  \end{columns}
  \mlsyscredit{Napkin Math 18.1, ``The delay wall''; stopping relation and positive root from The Physical Body (equations 2.2, 2.3). Values derived in images/calcs.py.}
  \note{[3.5 min] Worked example 1, the chapter's Napkin Math 18.1. ENGAGE: two minutes in pairs for the stop at 1.50 m/s. Step 2: the delay chain carries the carriage 96 mm before braking begins, already past the 60 mm stop line; braking adds 75 mm, a stop of 171 mm and a deficit of 111 mm. Step 3: the speed ceiling is the positive root of the stopping equation from Lecture 2, with the 45 mm left after the reserve: a times tau is 0.96 m/s, 2ad is 1.35, the root is 1.51, and the ceiling is 0.55 m/s (derived). The chain is 16 ms to a usable image, 36 ms for transfer, inference, scheduling, and admission, 2 ms on the bus, and 10 ms of brake buildup. WARN: students blame the brake. At these speeds delay travel exhausts the clearance before braking does; the speed limit belongs to the whole chain from closure to braking. FLEX: [CORE].}
\end{frame}

\begin{frame}{At 0.50 m/s the delay chain still sets most of the stop}
  \centering
  \mlsyssteps[width=0.92\textwidth,height=4.0cm,keepaspectratio]{cell_stop_step}{2}\\[0.15cm]
  \uncover<2->{\small Delay travel falls 3$\times$ with speed; braking falls 9$\times$. Kinetic energy falls from 13.5 J to 1.5 J.}
  \mlsyscredit{Slide variant of \S18.2.2, figure 2c. Derived in images/calcs.py; moving mass 12 kg (illustrative).}
  \note{[1 min] NARRATE: Step 1, the nominal stop: 96 mm of delay and 75 mm of braking against the 60 mm stop line, inside which a 15 mm reserve is held back. Step 2, the candidate mode at 0.50 m/s, under the 0.55 m/s ceiling: 32 mm of delay and 8.3 mm of braking, 40.3 mm, leaving 19.7 mm before the line. Delay is linear in speed and falls three{}fold; braking is quadratic and falls nine{}fold. The carriage's 12 kg carries 13.5 J at full speed and 1.5 J in the mode, 88.9 percent less (derived). ENGAGE: at 0.50 m/s, which term would you shorten next? Expected: the delay chain, since it is now four fifths of the stop. FLEX: [OPTIONAL].}
\end{frame}

\begin{frame}{A shorter chain buys speed; a weaker brake voids the mode}
  \begin{columns}[c]
    \begin{column}{0.62\textwidth}
      \centering\safeimg[width=\textwidth,height=5.6cm,keepaspectratio]{cell_ceilings.pdf}
    \end{column}
    \begin{column}{0.36\textwidth}
      \centering
      \safeimg[width=0.27\textwidth,height=4.4cm,keepaspectratio]{detec4-light-curtain.png}\hspace{0.12cm}%
      \raisebox{1.6cm}{\parbox{0.68\textwidth}{\small\raggedright\textbf{SICK deTec4}\\ safety light curtain\\[0.1cm]\footnotesize Response 10 ms (uncoded)\\ PL e, SIL 3\\ PFH$_\text{D}$ $15.3{\times}10^{-9}$ per h}}
    \end{column}
  \end{columns}
  \mlsyscredit{SICK deTec4 Set 30/450 product data sheet (2024), response time and safety parameters. Ceilings derived in images/calcs.py from Napkin Math 18.1 inputs; curtain chain 10 + 2 + 10 ms.}
  \note{[2.5 min] LINK: Lecture 12's certified floor beneath an unrated layer. NARRATE: the positive root shows how the ceiling responds to each premise. Halving the 36 ms processing allocation raises it from 0.55 to 0.66 m/s. Halving the qualified brake to 7.5 m/s squared lowers it to 0.47 m/s, under the 0.50 m/s mode, which then no longer fits. The two premises play different roles. Shorter latency is a throughput opportunity, while a braking shortfall invalidates the mode, so the loaded deceleration must be qualified before commissioning. Right: the real artifact. The camera-triggered stop is an application path; the cell's protective function lies elsewhere. A safety-rated light curtain such as this one feeds the drive's safe-stop input directly; its data sheet gives a 10 ms response, performance level e, and SIL 3. With command delivery and brake onset unchanged, the chain is 22 ms and the ceiling 0.88 m/s (derived; this slide's own arithmetic, which leaves out the drive's own reaction time). WARN: the 1 kHz monitor tick bounds only how long a command waits for the next check; it does not shorten transfer, inference, and scheduling. FLEX: [CORE].}
\end{frame}

\begin{frame}{The speed limit moves the proposer onto unevaluated states}
  \begin{columns}[c]
    \begin{column}{0.56\textwidth}
      \centering\safeimg[width=\textwidth,height=5.0cm,keepaspectratio]{derate_gap.pdf}
    \end{column}
    \begin{column}{0.41\textwidth}
      \small A slower carriage changes frame timing and motion blur.\\[0.2cm]
      Two ways to close the gap:\\[0.05cm]
      evaluate at the derated speed before release, or\\[0.05cm]
      record the regime as an open premise.
    \end{column}
  \end{columns}
  \mlsyscredit{Engineering Physical AI, \S18.2.2 (second law). Distribution shapes illustrative.}
  \note{[1.5 min] LINK: Lecture 5's endogenous covariate shift: a policy whose actions choose its next observations. NARRATE: the cell runs three learned proposers, for carriage speed from the contents, insertion force from carton thickness, and acceleration profiles from winding temperature. If closed-loop evaluation sampled only the nominal 1.50 m/s, clamping the carriage to 0.50 m/s puts the vision model in visual dynamics it never met. The restriction that secures the stop opens a regime without evidence. ENGAGE: which audit question catches this? Expected: step 2, run inside the envelope that step 1 set. WARN: students assume slower is always safer for the model; slower is safer for the stop and unevaluated for the proposer. FLEX: [CORE].}
\end{frame}

\begin{frame}{The cell earns a restricted mode and no autonomous sign-off}
  \begin{columns}[c]
    \begin{column}{0.64\textwidth}
      \footnotesize
      \renewcommand{\arraystretch}{1.3}
      \begin{tabular}{@{}P{0.30\textwidth}P{0.20\textwidth}P{0.38\textwidth}@{}}
        \toprule
        \textbf{Claim} & \textbf{Verdict} & \textbf{Premise} \\
        \midrule
        Transit at 0.50 m/s & \textcolor{nervecol}{\textbf{candidate}} & 64 ms, 15 m/s$^2$, 60 mm, reserve validated together \\
        Insertion force, thermal duty & \textcolor{delaycol}{\textbf{restricted}} & crush limits and sensor thermal lag untested \\
        Autonomous tab sign-off & \textcolor{conseq}{\textbf{refused}} & glare hides the tab; no independent sensor \\
        \bottomrule
      \end{tabular}
    \end{column}
    \begin{column}{0.33\textwidth}
      \centering\safeimg[width=\textwidth,height=3.6cm,keepaspectratio]{case-erector.jpg}\\[0.05cm]
      {\scriptsize A real case erector: carton forming and sealing}
    \end{column}
  \end{columns}
  \mlsyscredit{Engineering Physical AI, \S18.2.2 (fourth law). Photo: Killgore1138, CC BY-SA 4.0, Wikimedia Commons.}
  \note{[1.5 min] LINK: Lecture 16's signed verdicts. NARRATE: the fourth law turns the audit into a release argument. Transit at 0.50 m/s is a candidate, provided the 64 ms chain, the 15 m/s squared brake, the 60 mm clearance, and the reserve are validated together on the physical cell. Insertion forces and thermal duty stay restricted: a 220 N force limiter and a 105 degree thermistor trip are supervisory cutoffs that prove nothing about crush limits or sustained duty. Autonomous package sign-off is refused, because glare from tape and crinkled plastic both mimics and hides tab defects, and the cell has no tactile switch or laser profile scanner. ENGAGE: why does more glare-corrupted training data not fix the sign-off? Expected: the image lacks the information; it needs a new sensor channel. Mandatory manual inspection downstream supports a deployment that claims no autonomous sign-off. FLEX: [CORE].}
\end{frame}

\begin{frame}{Does 0.50 m/s survive a 12 m/s$^2$ brake?}
  \mlsysactivity{Your turn}{2 min}
  \begin{columns}[T]
    \begin{column}{0.42\textwidth}
      \begin{mlsyscard}{mlsysaccent}
        \small Load testing qualifies the brake at only \textbf{12 m/s$^2$}.\\
        Same 64 ms chain, same 45 mm budget.\\[0.1cm]
        \textbf{New ceiling? Does the mode fit?}
      \end{mlsyscard}
    \end{column}
    \begin{column}{0.54\textwidth}
      \small
      \uncover<2->{$a\tau = 0.77$ m/s;\enspace $2ad = 1.08$ m$^2$/s$^2$\\[0.1cm]
      Ceiling \textbf{0.52 m/s}\\[0.15cm]}
      \uncover<3->{At 0.50 m/s: 32 + 10.4 = \textbf{42.4 mm} of 45 mm.\\[0.1cm]
      It fits, with 2.6 mm of the budget left.}
    \end{column}
  \end{columns}
  \mlsyscredit{Checkpoint 18.1, first question. Derived in images/calcs.py.}
  \note{[2 min] Checkpoint 18.1 as a two-minute exercise. ENGAGE: pairs compute the ceiling. Step 2: a times tau is 12 times 0.064, 0.77 m/s; 2ad is 1.08; the positive root gives 0.52 m/s (derived). Step 3: at 0.50 m/s the delay travel stays 32 mm and braking grows to 10.4 mm, 42.4 mm of the 45 mm budget, so the mode still fits with 2.6 mm left inside the budget and the 15 mm reserve untouched. WARN: students who subtract from the 60 mm line instead of the 45 mm budget spend the reserve. The margin is now thin enough that the next premise to weaken, pad wear or a longer chain, would void the mode. Land segment 2: the audit transfers as questions and budgets, and each answer constrains the next. FLEX: [OPTIONAL].}
\end{frame}

% =============================================================================
% SECTION: UNOBSERVED PREMISES
% =============================================================================
\mlsysnote{[0.5 min] Segment 3 returns to the hook. What happens where the evidence runs out? FLEX: [CORE].}
\mlsysdivider{What can no sensor on the machine tell us?}{Engineering Physical AI, \S18.3}
\section{Unobserved}

\begin{frame}{Below $\mu$ = 0.204, floor friction sets the deceleration}
  \begin{columns}[c]
    \begin{column}{0.58\textwidth}
      \centering\safeimg[width=\textwidth,height=5.2cm,keepaspectratio]{friction_knee.pdf}
    \end{column}
    \begin{column}{0.39\textwidth}
      \normalsize $a = \min(a_{\text{brake}},\ \mu g)$\\[0.3cm]
      \small $\mu$: tire-floor friction coefficient\\
      $a_{\text{brake}}$ = 2 m/s$^2$\\[0.2cm]
      Floor binds when $\mu < 2/9.81 = 0.204$
    \end{column}
  \end{columns}
  \mlsyscredit{Engineering Physical AI, \S18.3, equation 1 and Napkin Math 18.2; friction values illustrative. Derived in images/calcs.py.}
  \note{[2 min] Worked example 2 begins. LINK: Lecture 2's credible deceleration, 2 m/s squared, was always conditional on the floor. NARRATE: the deceleration the machine actually has is the smaller of what the brakes deliver and what tire-floor friction transmits, mu times g, taken with all the weight on the driven wheels, the optimistic case. The knee sits at mu equal to 2 over 9.81, 0.204. Dry concrete at 0.60 is on the flat part; the brake binds. An inspected floor at 0.12 gives 1.18 m/s squared, and the film at 0.05 gives 0.49. ENGAGE: on the film, does a stronger brake help? Expected: no; to the left of the knee the floor binds, and more brake torque only locks the wheels. FLEX: [CORE].}
\end{frame}

\begin{frame}{Measured: a smooth tire on engine oil held $\mu$ = 0.01}
  \centering
  \safeimg[width=0.92\textwidth,height=5.0cm,keepaspectratio]{trl_friction.pdf}\\[0.05cm]
  \small Friction belongs to the tire, the surface, and what lies between them.
  \mlsyscredit{Measured: Lambourn and Viner, ``Friction Tests on Contaminated Road Surfaces,'' TRL PPR073, 2006, table 5.1 (locked-wheel averages; smooth tire = ASTM test tire). Floor line from equation 1.}
  \note{[2 min] Real data behind the illustrative 0.05. NARRATE: the UK Transport Research Laboratory locked the wheels of two skid cars and of a trailer with a smooth ASTM test tire on contaminated asphalt and concrete. Dry concrete gave 0.85 with treaded tires. Diesel on concrete gave 0.44 with treaded tires and 0.24 with the smooth one. Waste engine oil on asphalt gave 0.41 with treads and 0.01 with the smooth tire, far below the 0.204 floor. A robot's polyurethane wheel on polished concrete resembles the smooth tire on a low-texture surface more than a treaded car tire on a road. WARN: no one number is the friction of a floor; the coefficient belongs to the pair of surfaces and the film between them, which is why the braking premise needs a measurement at the interface, and no camera viewing the floor supplies it. FLEX: [CORE].}
\end{frame}

\begin{frame}{On the film the budget grows from 969.0 mm to 2.13 m}
  \centering
  \mlsyssteps[width=0.92\textwidth,height=4.3cm,keepaspectratio]{film_terms_step}{2}\\[0.1cm]
  \uncover<2->{\small Faster admission can recover only the 173.7 mm of delay travel.}
  \mlsyscredit{Napkin Math 18.2, ``The friction floor''; known-component budget at 1.3 m/s, excluding the dispatch-to-admission interval. Derived in images/calcs.py.}
  \note{[1 min] NARRATE: Step 1, the dry budget from Lecture 12 at 1.3 m/s: 173.7 mm of delay travel, 422.5 mm of braking, 182.8 mm for the jerk-limited ramp of Lecture 11, and 190 mm of fixed overheads, 969.0 mm against 1{,}100. Step 2, the same terms on the film: braking grows about four{}fold, to 1{,}723.3 mm, and the budget to 2.13 m (derived). Below the floor the braking term grows as one over mu. ENGAGE: what would a twice-as-fast enforcer buy? Expected: at most part of 173.7 mm, nowhere near the 1 m overrun. Only a restriction or a measurement at the tire-floor interface keeps the stop inside the clear distance. FLEX: [OPTIONAL].}
\end{frame}

\begin{frame}{Identical sensor histories leave no classifier a signal}
  \centering
  \safeimg[width=0.92\textwidth,height=4.1cm,keepaspectratio]{indist.pdf}\\[0.1cm]
  \small\raggedright\textbf{Observational indistinguishability:} two states that need different responses give the same observations, under every admissible input, until the deadline for a state-specific response.
  \mlsyscredit{Engineering Physical AI, \S18.3.4, definition 18.1 and equation 4. Traces illustrative.}
  \note{[2.5 min] The chapter's definition 18.1. NARRATE: two worlds, the dry floor and the film. Camera and lidar histories are identical until the base is already inside its dry-floor stopping distance. The deadline is set in distance: the remaining distance to the hazard must cover speed times the enforcer's and actuator's delay plus the braking distance, the same stopping relation as always. ENGAGE: will a better classifier on the same cameras find the film in time? Expected: no; identical histories give every classifier identical inputs. WARN: three distinctions. Noise and weak detectors hide information that more data can recover; here the information is absent from the channel. Lecture 9's unobservability asks whether a state can ever be reconstructed; indistinguishability is tied to a deadline and to the inputs the machine may apply. And a designed probe can sometimes make the states differ before the deadline, the subject of the next slides. FLEX: [CORE].}
\end{frame}

\begin{frame}{Which deadline, stop time or distance over speed?}
  \mlsysactivity{Vote}{1.5 min}
  \centering
  \only<1|handout:0>{\vspace{0.4cm}\normalsize An engineer sets the latest useful observation by subtracting braking time from the impact time.\\[0.3cm]
    \choice{A}{Subtract the stop's duration, 931.2 ms}\\[0.15cm]
    \choice{B}{Subtract braking distance over speed, 465.6 ms}}%
  \only<2-|handout:1>{\safeimg[width=0.92\textwidth,height=3.5cm,keepaspectratio]{overreserve.pdf}\\[0.15cm]
    \small \textbf{B.} The time form over-reserves 465.6 ms, or 605.3 mm, against a 100.8 ms spare.}
  \mlsyscredit{Engineering Physical AI, \S18.3.4 and \S18.5, pitfall ``Charging the full stop time against an action deadline.'' Derived in images/calcs.py.}
  \note{[1.5 min] ENGAGE: vote A or B. Step 2: B. The jerk-limited stop from 1.3 m/s lasts 931.2 ms but covers 605.3 mm, the distance of only 465.6 ms at constant speed, because the base slows as it brakes. Subtracting the full duration treats the machine as if it held speed and then stopped instantly. It over-reserves 465.6 ms, 605.3 mm, more than four times the 100.8 ms of timing allowance the whole budget leaves (derived). WARN: the error survives review because it is safe in the narrow sense; its cost is that it refuses modes the distance form correctly admits and hides how much allowance a new detector could still spend. FLEX: [OPTIONAL].}
\end{frame}

\begin{frame}{Five responses carry a claim when the state stays hidden}
  \centering
  \mlsyssteps[width=0.92\textwidth,height=5.0cm,keepaspectratio]{responses_step}{2}\\[0.1cm]
  \uncover<2->{\small The register already used two of them, a lower speed and no uncontrolled routes.}
  \mlsyscredit{Engineering Physical AI, \S18.3.4, ``Responses when the state remains hidden''; auxiliary signal design after Campbell and Nikoukhah 2004.}
  \note{[2 min] NARRATE: Step 1. First ask whether an admissible probe can separate the states before the deadline. On the floor, a brief torque pulse at low speed before the base enters a route makes the wheels slip measurably on a film; it turns friction under the tires into telemetry, but only for the patch the wheels stand on. At the latch, a tug at a qualified proof load checks that it holds that load. Step 2: if no probe works in time, four responses remain. A proof test measures from outside the runtime channel at intervals, such as an unloaded tare of the drifting wrist force sensor. A conservative action is safe in both states, such as a speed low enough for the film. Passive containment caps the consequence mechanically. A restricted domain withholds authority from states that cannot be told apart. ENGAGE: classify a series spring at the latch. Expected: passive containment. FLEX: [CORE].}
\end{frame}

\begin{frame}{A series spring keeps the latch under 100 N with torque cut}
  \begin{columns}[c]
    \begin{column}{0.60\textwidth}
      \centering\safeimg[width=\textwidth,height=4.9cm,keepaspectratio]{latch_bars.pdf}
    \end{column}
    \begin{column}{0.37\textwidth}
      \centering\safeimg[width=0.75\textwidth,height=3.0cm,keepaspectratio]{door-latch.jpg}\\[0.05cm]
      {\scriptsize A spring-free gate latch}\\[0.15cm]
      \raggedright\small Spring at most 50.5 kN/m, stroke at least 2.0 mm
    \end{column}
  \end{columns}
  \mlsyscredit{Engineering Physical AI, \S18.3.4, equation 5 and its margin figure; derived in images/calcs.py. Photo: Gaius Cornelius, CC BY-SA 3.0, Wikimedia Commons.}
  \note{[1.5 min] LINK: Lecture 12 owns passive-compliance stiffness; this chapter sizes it. NARRATE: if the safety response cuts motor torque at the 15 N tripwire, the arm's 12 kg effective mass coasts into the latch and energy sets the peak force. Against the rigid strike plate, 4.0 times ten to the five N/m, the guarded 0.03 m/s approach peaks at 67.4 N; the faster 0.10 m/s approach peaks at 219.6 N, far past the 100 N declared floor. With a series spring of at most 50.5 kN/m and 2.0 mm of stroke, the faster approach peaks at 74.9 N (all derived). WARN: the 100 N is a declared floor, a judgment recorded in the register; nobody has measured the latch's strength. Accountable judgment chooses the bound; it cannot create material strength, so proof tests must still check fatigue and wear. FLEX: [OPTIONAL].}
\end{frame}

\begin{frame}{The register binds each open premise to four things}
  \footnotesize
  \renewcommand{\arraystretch}{1.25}
  \begin{tabular}{@{}P{0.26\textwidth}P{0.37\textwidth}P{0.29\textwidth}@{}}
    \toprule
    \textbf{Missing observable} & \textbf{Restriction and owner} & \textbf{Closure test} \\
    \midrule
    Friction under a clear film & 1 m/s on inspected floors, $\mu \ge 0.12$; site operations & Film detector, 2{,}995 seeded trials, no miss \\
    \uncover<2->{A person walking out of the hidden cross aisle} & \uncover<2->{Barriers block cross traffic during runs; site operations} & \uncover<2->{Presence sensor qualified the same way} \\
    \uncover<3->{Object identity and crush strength} & \uncover<3->{Declared item list; singulation; site operations with the task profile} & \uncover<3->{Identity check, 2{,}995 seeded look-alikes} \\
    \bottomrule
  \end{tabular}\\[0.15cm]
  \small\uncover<3->{Each entry: an observable, a restriction, an owner, a closure test.}
  \mlsyscredit{Engineering Physical AI, \S18.3.3, table 2.}
  \note{[2.5 min] The lecture's own record, the residual-claims register, the sixteenth record in the chain. NARRATE: one entry per step. Step 1: the film. The premise is that every aisle floor supplies at least 0.204; the missing observable is friction under a clear film. The restriction is 1 m/s on routes whose floors are inspected to at least 0.12, under a known-component ceiling of 1.23 m/s; facility operations owns it through the spill procedure, because a software speed setting cannot keep a floor dry. The closure test is a film detector with zero misses in 2{,}995 seeded trials, a 0.1 percent miss bound at 95 percent confidence by Lecture 7's zero-failure sizing. Step 2: a person who keeps walking toward the base during a stop. Step 3: the glass jar grounded as the mug, which Lecture 10 showed passes every geometric check, since admission checks physics and cannot check meaning. WARN: a warning note in the documentation is not an entry. Until the closure test passes, the restriction stays wired into the signed limits and the site procedure, and every retrained model inherits it. Lecture 17's open human-role and detection premises enter in the same form. FLEX: [CORE].}
\end{frame}

\begin{frame}{What is the register entry for a worn latch?}
  \mlsysactivity{Diagnose}{2 min}
  \begin{columns}[c]
    \begin{column}{0.30\textwidth}
      \centering\safeimg[width=\textwidth,height=4.0cm,keepaspectratio]{door-latch.jpg}
    \end{column}
    \begin{column}{0.66\textwidth}
      \small Fatigue can lower the latch's yield force below the arm's contact ceiling while the joint torque sensors read the same.\\[0.15cm]
      \only<1|handout:0>{\normalsize Write its four parts.}%
      \only<2-|handout:1>{\footnotesize
        \textbf{Observable:} remaining yield margin of the latch\\
        \textbf{Restriction:} guarded 0.03 m/s approach; series spring\\
        \textbf{Owner:} facility maintenance, with the task profile\\
        \textbf{Closure:} proof tests at an interval that bounds wear}
    \end{column}
  \end{columns}
  \mlsyscredit{Checkpoint 18.2, second question; \S18.3.4. Photo: Gaius Cornelius, CC BY-SA 3.0, Wikimedia Commons.}
  \note{[2 min] ENGAGE: pairs write the four parts. Step 2: one acceptable answer. The missing observable is the latch's remaining yield margin, which the torque channel does not carry. The restriction holds the guarded approach and fits passive compliance, so force stays under the declared floor. An owner who can keep the premise true, such as facility maintenance with the task profile. The closure test is a scheduled proof load, at an interval short enough that maintenance records show wear cannot carry the yield force below the ceiling between tests. WARN: students write ``retrain the contact model'' as the closure test; no model on the torque channel can see a margin the channel does not carry. Land segment 3: where no sensor reports a premise, the machine gets a restriction, an owner, and a test, decided before operation. FLEX: [CORE].}
\end{frame}

% =============================================================================
% SECTION: NEXT CAPABILITY
% =============================================================================
\mlsysnote{[0.5 min] Segment 4 prices what a shield buys and asks which experiment the next effort should fund. FLEX: [CORE].}
\mlsysdivider{Which experiment should the next effort buy?}{Engineering Physical AI, \S18.3.2 and \S18.4}
\section{Next Capability}

\begin{frame}{A shield lowers a hazard rate only by its tested coverage}
  \centering
  \vspace{0.3cm}
  {\Large $\lambda_{\text{sys}} \le \underbrace{\lambda_h}_{\text{hazardous proposals}}\,\big(1 - \underbrace{c_h}_{\text{coverage}}\big) + \underbrace{\lambda_{\text{cc}}}_{\text{common cause}}$}\\[0.5cm]
  \small\raggedright
  \uncover<2->{Rates per operating hour; $c_h$ = probability the shield intercepts a hazardous proposal of class $h$.\\[0.15cm]}
  \uncover<3->{\textcolor{conseq}{Fleet rates already include the shield, so multiplying them by $(1 - c_h)$ counts it twice.}\\[0.1cm]
  \textcolor{conseq}{A film the shield cannot sense earns no credit.}}
  \mlsyscredit{Engineering Physical AI, \S18.3.2, equation 2 (shield dilution).}
  \note{[2 min] The chapter's second-order concept, shield dilution. NARRATE: Step 1, the terms. Lambda h is how often the policy emits hazardous proposals before the permission layer. c h is the shield's tested interception coverage for that hazard class. Lambda cc is hardware and common-cause failure. With perfect coverage only the common-cause term remains; with none, every hazardous proposal reaches the actuators. Step 2: units. Step 3: two traps. A fleet's observed incident rate already includes every protective layer, so multiplying it by an interception factor double-counts. And the equation gives no credit for hazards the shield cannot sense, such as the film or a fatigued latch; those stay in the register. WARN: an occluded frame that makes the policy propose a bad trajectory can blind the shield's estimator at the same moment, lowering c h exactly where the policy errs. Coverage must be measured on the states where the policy fails. FLEX: [CORE].}
\end{frame}

\begin{frame}{Does a shield make proposer errors irrelevant?}
  \mlsysactivity{Vote}{2.5 min}
  \begin{columns}[T]
    \begin{column}{0.42\textwidth}
      \begin{mlsyscard}{mlsysaccent}
        \small\textbf{Given} \smalltag{(illustrative)}\\[0.05cm]
        100 machines, 90 days, 24 h: zero hazardous proposals\\
        Coverage $c_h$ = 0.9999; $\lambda_{\text{cc}}$ = 0\\
        Target $10^{-9}$ per hour\\[0.05cm]
        \textbf{Is the target met?}
      \end{mlsyscard}
    \end{column}
    \begin{column}{0.54\textwidth}
      \small
      \uncover<2->{Proposal bound: $2.996 / 2.16{\times}10^{5}$ h = $1.387{\times}10^{-5}$ per h\\
      Residual: $1.387{\times}10^{-9}$ per h \textcolor{conseq}{\textbf{above target}}\\[0.2cm]}
      \uncover<3->{Coverage itself needs \textbf{29{,}956} seeded hazards with no miss.\\
      Meeting the bound needs 124.8 days of log, 34.8 beyond the cycle.}
    \end{column}
  \end{columns}
  \mlsyscredit{Napkin Math 18.3, ``The coverage bill''; rule of three (Hanley and Lippman-Hand 1983). Derived in images/calcs.py.}
  \note{[2.5 min] Worked example 3 and the chapter's fallacy. ENGAGE: vote yes or no, then a minute with a neighbor. Step 2: no. Zero hazardous proposals in 216{,}000 proposal-hours bound the rate at minus ln 0.05, 2.996, over the hours, 2.996 divided by 216{,}000: 1.387 times ten to the minus five per hour, Lecture 7's rule of three with hours in place of trials. Times one minus 0.9999, the residual is 1.387 times ten to the minus nine, above the target even with the common-cause term set to zero. Step 3: the evidence the claim needs. Showing coverage of 0.9999 takes 29{,}956 seeded hazardous proposals with no miss; bringing the proposal bound down to ten to the minus five takes 124.8 days of log, 34.8 beyond the cycle (derived). WARN: the target is borrowed from civil aviation's extremely improbable band, a per-condition share of risk that no robotics standard requires. The shield lets the proposer be imperfect; the team must still measure how often it errs. FLEX: [CORE].}
\end{frame}

\begin{frame}{Measured fleet rates already include every protective layer}
  \begin{columns}[c]
    \begin{column}{0.36\textwidth}
      \centering\safeimg[width=\textwidth,height=3.4cm,keepaspectratio]{waymo-jaguar-sf.jpg}\\[0.05cm]
      {\scriptsize A driverless Waymo vehicle, San Francisco}
    \end{column}
    \begin{column}{0.60\textwidth}
      \small Lecture 16's Waymo data showed 183 injury crashes in 271.3 million rider-only miles, 0.67 per million against a benchmark of 3.77.\\[0.3cm]
      That rate already counts every layer the cars run.\\[0.3cm]
      \textcolor{conseq}{A team that multiplies it by its own shield's $(1 - c_h)$ counts the protection twice.}
    \end{column}
  \end{columns}
  \mlsyscredit{Measured: Waymo Safety Impact data hub (data cuts to June 2026), via MLSysIM, as shown in Lecture 16. Engineering Physical AI, \S18.3.2. Photo: Dllu, CC BY 4.0, Wikimedia Commons.}
  \note{[2 min] Real fleet data, recalled from Lecture 16 in one line: 183 injury-reported crashes in 271.3 million rider-only miles, 0.67 per million miles against a location-adjusted benchmark of 3.77. Today reads it with equation 2 in mind. ENGAGE: a team wants to claim its own shield would cut this rate another hundredfold. What is wrong? Expected: the measured rate already includes every protective layer the vehicles run, so multiplying it by an interception factor counts protection twice. NARRATE: the claim a team may make credits a new shield only with its tested coverage of the hazards the existing layers miss, measured with those layers in place; that coverage is what the seeded-proposal campaign of the previous slide prices. WARN: even at this scale, Lecture 7's exposure wall holds: hundreds of millions of miles bound injury rates well, and they do not reach the rates per hour that a borrowed 10 to the minus 9 target would demand. FLEX: [CORE].}
\end{frame}

\begin{frame}{The kind of missing evidence selects the next experiment}
  \begin{columns}[T]
    \begin{column}{0.49\textwidth}
      \begin{mlsyscard}[fontupper=\footnotesize\raggedright]{conseq}
        \gapcell{thumb_indist.pdf}{Unmeasured observable}{A new sensor channel with a tested miss bound.}
      \end{mlsyscard}
      \uncover<2->{\begin{mlsyscard}[fontupper=\footnotesize\raggedright]{brainnavy}
        \gapcell{thumb_friction_knee.pdf}{Theoretical deficit}{A bound whose premises are tested on the plant.}
      \end{mlsyscard}}
    \end{column}
    \begin{column}{0.49\textwidth}
      \begin{mlsyscard}[fontupper=\footnotesize\raggedright]{delaycol}
        \gapcell{thumb_derate_gap.pdf}{Evaluation deficit}{Field trials on the real population of states.}
      \end{mlsyscard}
      \uncover<2->{\begin{mlsyscard}[fontupper=\footnotesize\raggedright]{plantcol}
        \gapcell{unitree-g1.jpg}{Scope mismatch}{A new architecture or a narrower domain.}
      \end{mlsyscard}}
    \end{column}
  \end{columns}
  \mlsyscredit{Engineering Physical AI, \S18.4.1. Thumbnails of earlier figures in this lecture. Photo: Sayanesy, CC0, Wikimedia Commons.}
  \note{[2 min] The third concept the lecture owns: next capability by missing evidence. Each card carries the shape of the figure from this lecture that showed that kind of gap: the identical sensor histories, the states admitted against the states evaluated, the friction bound, and a balancing humanoid. NARRATE: Step 1, two classes. An unmeasured observable needs a sensor that separates the states before the deadline: the film, the hidden cross-aisle person, object identity, the packaging tab. An evaluation deficit needs trials or estimators that represent the real population, such as the shield coverage of the last slide. Step 2, two more. A theoretical deficit needs an analytical bound whose physical premises can be proved or tested. A scope mismatch cannot be fixed by any measurement, trial, or proof inside the current architecture: cutting torque on a balancing biped is itself the hazard. ENGAGE: classify ``loaded dispatch-to-admission timing,'' the allocation of Lecture 13. Expected: an evaluation deficit, closed by timing analysis challenged with loaded and adversarial measurements. WARN: training another neural detector closes none of these when the evidence missing is a field evaluation, a bound, or a fallback the body cannot execute. Moving an unverified friction assumption into a simulator leaves it as unverified as before. FLEX: [CORE].}
\end{frame}

\begin{frame}{A 23 percent speed cut costs 30 percent more time per leg}
  \begin{columns}[c]
    \begin{column}{0.50\textwidth}
      \centering
      {\Huge\bfseries\textcolor{mlsysaccent}{+30\%}}\\[0.1cm]
      {\small 30 m leg: 23.1 s at 1.3 m/s, 30 s at 1 m/s}\\[0.4cm]
      {\Large\bfseries\textcolor{plantcol}{2{,}995}}\\[0.05cm]
      {\small seeded-film trials to retire the restriction}
    \end{column}
    \begin{column}{0.46\textwidth}
      \small A restriction stays until its closure test passes.\\[0.2cm]
      The campaign pays back when recovered throughput exceeds the test and sensor cost.\\[0.2cm]
      Cheap upkeep of the premise, such as a spill procedure on a few routes, can win.
    \end{column}
  \end{columns}
  \mlsyscredit{Engineering Physical AI, \S18.4.2; derived in images/calcs.py.}
  \note{[1.5 min] NARRATE: the film entry costs throughput. Lowering cruise from 1.3 to 1 m/s is a 23.1 percent speed cut, but transit time scales as one over speed, so a 30 m leg grows from 23.1 to 30 s, 30 percent longer (derived), before acceleration and docking. Retiring the restriction costs a campaign: 2{,}995 seeded-film trials plus the instrument. The coverage campaign for the shield, 29{,}956 seeded proposals, does not close the film entry; each entry has its own test. ENGAGE: when is keeping the restriction the right call? Expected: when the premise is cheap to keep true, such as cleaning and spill control on a few pallet-jack routes. Commercial fleets bill on work done: Agility Robotics reported 100{,}000 totes moved at a first customer site by late 2025, so a withheld mode is an operating expense. WARN: Bainbridge's ironies of automation: rare interventions leave operators unpracticed, so intervention rates carry a cost both ways. Lecture 17's worker burden belongs in the same ledger. FLEX: [CORE].}
\end{frame}

\begin{frame}{Six open problems, each tied to a law and a gap class}
  \centering\small
  \renewcommand{\arraystretch}{1.3}
  \begin{tabular}{@{}P{0.37\textwidth}P{0.33\textwidth}P{0.24\textwidth}@{}}
    \toprule
    \textbf{Open problem} & \textbf{Law stressed} & \textbf{Gap class} \\
    \midrule
    Cross-embodiment policies & permission must be earned & evaluation \\
    Data scaling & endogenous data & evaluation \\
    World models as evaluators & permission must be earned & evaluation \\
    Fall-safe balancing machines & irreversibility & scope mismatch \\
    Contact-rich manipulation & proposal is not permission & unmeasured observable \\
    Certifying learned safety functions & permission must be earned & theoretical \\
    \bottomrule
  \end{tabular}
  \mlsyscredit{Engineering Physical AI, \S18.4.3, table 3 (state of practice as of 2026).}
  \note{[2 min] NARRATE: the chapter's research frontier, each problem placed by the law it stresses and the evidence it lacks. Cross-embodiment policies such as pi zero need evidence that transfers across bodies and updates under a stated equivalence condition. Data scaling needs curves of closed-loop failure rates inside a declared ODD in place of average success. World models such as Cosmos need a proxy-transfer condition validated on the failure tail before their rollouts count as release evidence. A balancing humanoid like the one on the last slide has no resting state to fall into when torque is cut; certifying an actively controlled reaction to a supported rest is a scope mismatch with fail-passive safety standards. Contact needs a qualified tactile channel. Certifying a learned safety function needs a composite argument of proofs and statistics; the EU requires third-party assessment of such components from 20 January 2027. ENGAGE: pick one and name the first experiment. Land segment 4: the kind of evidence an open premise lacks decides the work. FLEX: [OPTIONAL].}
\end{frame}

% =============================================================================
% SECTION: SYNTHESIS
% =============================================================================
\mlsysnote{[0.5 min] The synthesis returns to Lecture 0's aisle run and to the hook. FLEX: [CORE].}
\mlsysdivider{What may the warehouse machine be trusted to do?}{Engineering Physical AI, \S18.5 and \S18.6}
\section{Synthesis}

\begin{frame}{Eight lectures built the aisle budget; one film breaks it}
  \centering
  \only<1|handout:0>{\mlsyskitfig[width=0.95\textwidth,height=4.8cm,keepaspectratio]{budget_bar_aisle_step7}}%
  \only<2-|handout:1>{\mlsyskitfig[width=0.95\textwidth,height=4.8cm,keepaspectratio]{budget_bar_aisle_step8}}\\[0.05cm]
  \only<1|handout:0>{\small The known components at 1.3~m/s use 969.0~mm of 1{,}100.}%
  \only<2|handout:0>{\small The 63.3~ms dispatch-to-admission allocation leaves 48.7~mm.}%
  \only<3-|handout:1>{\small On the film the same rows reach \textbf{2.13~m}. Faster software cannot close the gap.}
  \mlsyscredit{Aisle run, Volume IV course plan; stopping-budget bar from the course kit (MLSysIM \texttt{WarehouseAisle}); film value from images/calcs.py. Illustrative and derived values.}
  \note{[2.5 min] Lecture 0 asked how far a robot travels while it decides. This is the finished answer. NARRATE: Step 1: the known components at the 1.3 m/s aisle speed. Delay travel, 173.7 mm, from the elementary budget of Lecture 1, the chunk lease of Lecture 4, and observation age of Lecture 8; braking, 422.5 mm, with Lecture 2's credible deceleration; the jerk ramp of Lecture 11, 182.8 mm; and Lecture 12's fixed overheads, 190 mm. Step 2: Lecture 13's 63.3 ms dispatch-to-admission allocation, 82.3 mm, leaves the 48.7 mm spare Lecture 16 released against (derived). Step 3: today's row. The film multiplies the braking term and the stop reaches 2.13 m. ENGAGE: what answers the hook? Expected: a restriction with an owner and a closure test, or a measurement at the tire-floor interface; no faster model or enforcer does. FLEX: [CORE].}
\end{frame}

\begin{frame}{Sixteen records carry the argument from body to release}
  \centering
  \only<1|handout:0>{\mlsyskitfig[width=0.92\textwidth,height=5.2cm,keepaspectratio]{record_chain_build1}}%
  \only<2|handout:0>{\mlsyskitfig[width=0.92\textwidth,height=5.2cm,keepaspectratio]{record_chain_build2}}%
  \only<3|handout:0>{\mlsyskitfig[width=0.92\textwidth,height=5.2cm,keepaspectratio]{record_chain_build3}}%
  \only<4-|handout:1>{\mlsyskitfig[width=0.92\textwidth,height=5.2cm,keepaspectratio]{record_chain_l18}}
  \mlsyscredit{The record chain, The Causal Boundary, \S1.11; residual-claims register, \S18.3.3.}
  \note{[1 min] NARRATE: Lecture 1's record chain, rebuilt one part per step and now complete: one record per lecture from the limit record of Lecture 2 to the approval record of Lecture 17, closed by the residual-claims register. Step 1: Part I, the limit and handoff records. Step 2: Part II, provenance, policy manifest, and evaluation. Step 3: Part III, observation to placement, where the two strands meet at enforcement. Step 4: Part IV, authority to approval, and today's register, lit. Lecture 3 adds no record of its own; its proposer limits enter the handoff record. Read the chain as one argument, in which each record is the evidence the next decision reads. The register is the only one that records what the evidence does not support. FLEX: [CORE].}
\end{frame}

\begin{frame}{Five ideas to carry out of this course}
  \only<1|handout:0>{\mlsysactivity{Retrieval}{2 min}%
    \centering
    \vspace{0.7cm}
    {\Large\bfseries Question of the day}\\[0.5cm]
    {\large Answer the course's question in three sentences:\\[0.15cm] why does a better model not make a better machine?}\\[0.5cm]
    {\normalsize\textcolor{midgray}{Two minutes. Hand it in on the way out.}}}%
  \only<2-|handout:1>{%
    \small
    \thumbline{thumb_closed_loop_step2.pdf}{The four laws form one argument; a better prediction strengthens one link.}\\[0.12cm]
    \thumbline{thumb_audit_step4.pdf}{The audit transfers as four ordered questions; the old machine's numbers stay behind.}\\[0.12cm]
    \thumbline{thumb_film_terms_step2.pdf}{A premise no sensor reports gets a restriction, an owner, and a closure test.}\\[0.12cm]
    \thumbline{thumb_gaps.pdf}{The kind of missing evidence chooses the next experiment.}\\[0.12cm]
    \thumbline{thumb_budget_bar_aisle_step8.pdf}{The aisle budget is 969.0 mm dry and 2.13 m on the film, of 1{,}100 mm clear.}}
  \note{[2 min] Step 1: the course's exit ticket. Two minutes, from memory. A good answer names a premise the model cannot supply, such as the floor's friction, the states its own restriction creates, or a detector that does not exist, and says what the rest of the machine supplies instead. Collect the cards. Step 2: read the five lines, one per segment and the synthesis. If time is short, set the exit ticket as the last homework. FLEX: [CORE].}
\end{frame}

\begin{frame}{Open questions to carry beyond the course}
  \begin{columns}[c]
    \begin{column}{0.74\textwidth}
      \normalsize
      \begin{itemize}\setlength{\itemsep}{0.35cm}
        \item When a model updates every week from fleet data, which register entries must reopen, and who decides?
        \item Who should own a premise that lives in the building?
        \item When could a learned world model's rollouts count as release evidence?
      \end{itemize}
    \end{column}
    \begin{column}{0.23\textwidth}
      \centering\safeimg[width=\textwidth,height=4.0cm,keepaspectratio]{unitree-g1.jpg}
    \end{column}
  \end{columns}
  \mlsyscredit{Photo: Sayanesy, CC0, Wikimedia Commons.}
  \note{[1 min] No single answer. Ask students to write one sentence on any of the three. The first is continual learning against the register; the second is the film's owner, facility operations, generalized; the third is the world-model frontier. FLEX: [OPTIONAL].}
\end{frame}

\begin{frame}{Read, practice, and go further}
  \small
  \begin{itemize}\setlength{\itemsep}{0.12cm}
    \item \textbf{Read:} Engineering Physical AI, \S18.1 to \S18.6; Checkpoints 18.1 and 18.2.
    \item Butler and Finelli, ``The Infeasibility of Quantifying the Reliability of Life-Critical Real-Time Software,'' IEEE TSE, 1993.
    \item Sha, ``Using Simplicity to Control Complexity,'' IEEE Software, 2001.
    \item Hsu, Hu, and Fisac, ``The Safety Filter,'' Annual Review of Control, Robotics, and Autonomous Systems, 2024.
    \item Leveson, \emph{Engineering a Safer World}, MIT Press, 2011.
    \item Koopman, \emph{How Safe Is Safe Enough?}, 2022.
  \end{itemize}
  \note{[0.5 min] The chapter first, with both checkpoints; Checkpoint 18.1 applied to a machine students know reviews the whole course. Butler and Finelli show in a few pages why testing cannot validate ultra-high reliability. Sha gives the runtime-assurance argument behind the proposal and permission split. Hsu and colleagues survey safety filters and their coverage assumptions. Leveson treats safety as a control problem across the whole organization, the basis for owners on register entries. Koopman shows how explicit safety cases replace mileage claims. FLEX: [OPTIONAL].}
\end{frame}

\begin{frame}[c]
  \centering
  {\small\textcolor{lightgray}{FOR FURTHER STUDY}}\\[0.4cm]
  {\LARGE\bfseries Which premise of the next machine\\will no sensor report?}\\[0.6cm]
  \safeimg[width=0.30\textwidth,height=2.2cm,keepaspectratio]{erflex-mobile-manipulator.jpg}\\[0.3cm]
  {\footnotesize\textcolor{midgray}{Engineering Physical AI, \S18.6 and ``What's Next''}}
  \mlsyscredit{Photo: Auledas, CC BY 4.0, Wikimedia Commons.}
  \note{[0.5 min] There is no next lecture. The next unfamiliar machine brings a different body, task, and environment, and its engineering begins by tracing how observations become proposals and admitted commands become force, motion, and heat. Leave students with the question, then end. FLEX: [CORE].}
\end{frame}

\end{document}
